%%%%%%%%%%%%%%%%%%%%%%%%%%%%%%%%%%%%%%%%%
% Daily Laboratory Book
% LaTeX Template
%
% This template has been downloaded from:
% http://www.latextemplates.com
%
% Original author:
% Frank Kuster (http://www.ctan.org/tex-archive/macros/latex/contrib/labbook/)
%
% Important note:
% This template requires the labbook.cls file to be in the same directory as the
% .tex file. The labbook.cls file provides the necessary structure to create the
% lab book.
%
% The \lipsum[#] commands throughout this template generate dummy text
% to fill the template out. These commands should all be removed when
% writing lab book content.
%
% HOW TO USE THIS TEMPLATE
% Each day in the lab consists of three main things:
%
% 1. LABDAY: The first thing to put is the \labday{} command with a date in
% curly brackets, this will make a new page and put the date in big letters
% at the top.
%
% 2. EXPERIMENT: Next you need to specify what experiment(s) you are
% working on with an \experiment{} command with the experiment shorthand
% in the curly brackets. The experiment shorthand is defined in the
% 'DEFINITION OF EXPERIMENTS' section below, this means you can
% say \experiment{pcr} and the actual text written to the PDF will be what
% you set the 'pcr' experiment to be. If the experiment is a one off, you can
% just write it in the bracket without creating a shorthand. Note: if you don't
% want to have an experiment, just leave this out and it won't be printed.
%
% 3. CONTENT: Following the experiment is the content, i.e. what progress
% you made on the experiment that day.
%
%%%%%%%%%%%%%%%%%%%%%%%%%%%%%%%%%%%%%%%%%

%--------------------------------------------------------------------------------------
%	PACKAGES AND OTHER DOCUMENT CONFIGURATIONS
%--------------------------------------------------------------------------------------

\documentclass[idxtotoc,hyperref,openany,oneside]{ROLatex/labbook} % 'openany' here removes the gap page between days, erase it to restore this gap; 'oneside' can also be added to remove the shift that odd pages have to the right for easier reading

\newcommand{\HRule}{\rule{\linewidth}{0.5mm}} % Command to make the lines in the title page
\setlength\parindent{0pt} % Removes all indentation from paragraphs
\setlength{\parskip}{10pt plus 1pt minus 1pt} % Set space between paragraphs

% My packages
\usepackage{ROLatex/ROPackage} % This is in $TEXMFHOME/tex/latex
%--------------------------------------------------------------------------------------

\begin{document}

%--------------------------------------------------------------------------------------
%	TITLE PAGE
%--------------------------------------------------------------------------------------

\frontmatter % Use Roman numerals for page numbers
\title{
\begin{center}
\HRule \\[0.4cm]
{\Huge \bfseries Laboratory Journal \\[0.5cm] \Large Georgia Tech Robotics P.h.D}\\[0.4cm] % Degree
\HRule \\[1.5cm]
\end{center}
}
\author{\Huge Rowland O'Flaherty \\ \\ \LARGE rowlandoflaherty.com \\[2cm]}
\date{Beginning DD Month YYYY} % Beginning date
\maketitle

\tableofcontents

\mainmatter % Use Arabic numerals for page numbers

%--------------------------------------------------------------------------------------
%	LAB BOOK CONTENTS
%--------------------------------------------------------------------------------------

% Blank template to use for new days:

\labday{Dec 30, 2000}

\experiment{Experiment 1}

Some text for experiment 1 is this working

\begin{figure}[H]
  \centering
  \includegraphics[]{Figs_2000_12_30/example_figure.png}
  \caption{Caption here}
  \label{fig:1}
\end{figure}


\experiment{Experiment 2}

Some text for experiment 2


%--------------------------------------------------------------------------------------

\end{document}
